\section{从对称性推导振动与光谱}
\label{sec:vibrational-applications}
\learning{把原子位移约化为简正振动，分别数清振动自由度、不同基频和光谱允许的基频；能从特征标表读出红外与拉曼活性。}
\sources{D，第12--38页；谐振子、非谐性与转动见第14--20页。}

\subsection{先分清：模式数、频率数与谱带数}

一个含 $N$ 个原子的分子有 $3N$ 个笛卡儿位移坐标。整体平移占三个自由度；非线性分子的整体转动另占三个，因此振动自由度为 $3N-6$。线性分子绕自身轴的转动不改变原子核位置，独立转动只有两个，振动自由度为 $3N-5$。

简正模式是各原子按固定相对振幅、以同一频率协同运动的独立振动。若
\[
\Gamma_{\mathrm{vib}}=\sum_i a_i\Gamma_i,
\qquad n_{\mathrm{vib}}=\sum_i a_i d_i,
\]
其中 $a_i$ 是不可约表示 $\Gamma_i$ 出现的次数，$d_i$ 是其维数。一个 $E$ 表示贡献两个简正模式，却只给出一个二重简并频率；一个 $T$ 表示贡献三个模式和一个三重简并频率。相反，$2A_1$ 是两个独立的非简并振动，通常有两个不同频率，而不是“一个频率强度加倍”。

\begin{keyidea}[数峰时的三个层次]
先以维数检查模式总数，再把每个简并组算作一个频率，最后只保留光谱允许的组。例如 $2A_1+E$ 有四个模式、通常三个不同基频；若只有 $E$ 红外活性，则只有一个对称允许的红外基频。这里的理论谱带数按理想对称性、基频和无偶然重合计算；实际谱图还会受强度、分辨率、频率重合及对称性降低等因素影响。
\end{keyidea}

\subsection{全位移空间：不动原子法}

不必先猜出每个振动的箭头图。以全部原子的 $x,y,z$ 位移为基，先求 $3N$ 维表示 $\Gamma_{3N}$。对称操作 $R$ 若把某个原子送到另一原子的位置，其位移出现在非对角矩阵块中，对迹没有贡献；只有位置不变的原子贡献一个三维坐标变换矩阵的迹。所以
\[
\chi_{3N}(R)=n_{\mathrm{fix}}(R)\,\chi_{xyz}(R).
\]
$n_{\mathrm{fix}}$ 是操作后仍在原位置的原子数，$\chi_{xyz}$ 是极向量 $(x,y,z)$ 的变换矩阵之迹。常用值为：恒等操作 $3$、绕任意轴旋转角 $\theta$ 时 $1+2\cos\theta$、镜面反射 $1$、反演 $-3$。这些数也可由特征标表中 $x,y,z$ 所属表示的特征标相加得到。

把 $\Gamma_{3N}$ 按第\ref{sec:reduction}节的公式约化后，扣去平移与转动：
\[
\Gamma_{\mathrm{vib}}=\Gamma_{3N}-\Gamma_{\mathrm{trans}}-\Gamma_{\mathrm{rot}}.
\]
平移查表中的 $x,y,z$，转动查 $R_x,R_y,R_z$；线性分子只扣两个独立转动。最后必须用 $\chi_{\mathrm{vib}}(E)$ 检查是否等于 $3N-6$ 或 $3N-5$。

\begin{worked}[水：从九个位移到三个振动]
把 \ce{H2O} 放在 $yz$ 平面，O 位于原点，$z$ 轴为键角平分线和 $C_2$ 主轴，两颗 H 的坐标可写成 $(0,\pm a,b)$，其中 $a,b$ 为非零几何参数。严格采用下列列序：
\begin{center}
\small
\begin{tabular}{lrrrr}
\toprule
操作 & $E$ & $C_2(z)$ & $\sigma(xz)$ & $\sigma(yz)$\\
\midrule
$n_{\mathrm{fix}}$ & 3 & 1 & 1 & 3\\
$\chi_{xyz}$ & 3 & $-1$ & 1 & 1\\
$\chi_{3N}$ & 9 & $-1$ & 1 & 3\\
\bottomrule
\end{tabular}
\end{center}
例如 $C_2(z)$ 交换两个 H，只固定 O；O 的位移变成 $(-x,-y,z)$，故贡献 $-1-1+1=-1$。分子平面 $\sigma(yz)$ 固定全部三个原子，每个贡献 $-1+1+1=1$。

$C_{2v}$ 的群阶为 4，四个操作各自成类。逐行约化得到
\begin{align*}
a_{A_1}&=(9-1+1+3)/4=3, & a_{A_2}&=(9-1-1-3)/4=1,\\
a_{B_1}&=(9+1+1-3)/4=2, & a_{B_2}&=(9+1-1+3)/4=3.
\end{align*}
因而
\begin{align*}
\Gamma_{3N}&=3A_1+A_2+2B_1+3B_2,\\
\Gamma_{\mathrm{trans}}&=A_1+B_1+B_2,\\
\Gamma_{\mathrm{rot}}&=A_2+B_1+B_2,\\
\boxed{\Gamma_{\mathrm{vib}}=2A_1+B_2}.&
\end{align*}
三个模式分别是 $A_1$ 对称伸缩、$A_1$ 弯曲、$B_2$ 反对称伸缩。三者均非简并，通常有三个基频。由于 $z\sim A_1$、$y\sim B_2$，三个基频均红外允许；$A_1$ 含平方项，$B_2$ 含 $yz$，故也都拉曼允许。两个 $A_1$ 振动的红外跃迁偶极矩沿 $z$，$B_2$ 的沿 $y$。换坐标后标签可能变化，但不能只换 $B_1/B_2$ 而不同时换表和位移方向。
\end{worked}

\subsection{含简并的例子：平面三角形 \ce{BCl3}}

取分子平面为 $xy$，$C_3$ 主轴为 $z$，点群是 $D_{3h}$。下表给出本题需要的完整特征标，类名前的数字是该类操作数；撇号表示对分子平面反射为对称，双撇号表示反对称。
\begin{center}
\small
\begin{tabular}{lrrrrrrl}
\toprule
$D_{3h}$ & $E$ & $2C_3$ & $3C_2'$ & $\sigma_h$ & $2S_3$ & $3\sigma_v$ & 常用基\\
\midrule
$A_1'$ & 1 & 1 & 1 & 1 & 1 & 1 & $x^2+y^2,z^2$\\
$A_2'$ & 1 & 1 & $-1$ & 1 & 1 & $-1$ & $R_z$\\
$E'$ & 2 & $-1$ & 0 & 2 & $-1$ & 0 & $(x,y)$\\
$A_1''$ & 1 & 1 & 1 & $-1$ & $-1$ & $-1$ & ---\\
$A_2''$ & 1 & 1 & $-1$ & $-1$ & $-1$ & 1 & $z$\\
$E''$ & 2 & $-1$ & 0 & $-2$ & 1 & 0 & $(R_x,R_y)$\\
\bottomrule
\end{tabular}
\end{center}
二次函数另有 $(x^2-y^2,xy)\sim E'$、$(xz,yz)\sim E''$。注意 $E'$ 和 $E''$ 都是二维表示，不能把表中的一行只算成一个振动自由度。

\begin{worked}[\ce{BCl3}：六个模式为何只有四个基频]
沿用上表列序，不动原子法给出
\begin{center}
\small
\begin{tabular}{lrrrrrr}
\toprule
 & $E$ & $2C_3$ & $3C_2'$ & $\sigma_h$ & $2S_3$ & $3\sigma_v$\\
\midrule
$n_{\mathrm{fix}}$ & 4 & 1 & 2 & 4 & 1 & 2\\
$\chi_{xyz}$ & 3 & 0 & $-1$ & 1 & $-2$ & 1\\
$\chi_{3N}$ & 12 & 0 & $-2$ & 4 & $-2$ & 2\\
\bottomrule
\end{tabular}
\end{center}
群阶为 12。例如
\[
a_{A_1'}=\frac{12+2(0)+3(-2)+4+2(-2)+3(2)}{12}=1,
\]
而
\[
a_{E'}=\frac{12(2)+2(0)(-1)+3(-2)(0)+4(2)+2(-2)(-1)}{12}=3.
\]
其余行同法，得到
\begin{align*}
\Gamma_{3N}&=A_1'+A_2'+3E'+2A_2''+E'',\\
\Gamma_{\mathrm{trans}}&=E'+A_2'',\qquad
\Gamma_{\mathrm{rot}}=A_2'+E'',\\
\boxed{\Gamma_{\mathrm{vib}}=A_1'+2E'+A_2''}.&
\end{align*}
维数为 $1+2\times2+1=6=3\times4-6$。$A_1'$ 是对称伸缩，两个 $E'$ 分别对应一组二重简并伸缩和一组二重简并面内弯曲，$A_2''$ 是面外弯曲。通常共有四个不同基频，而不是六个。

红外查 $E'\sim(x,y)$ 和 $A_2''\sim z$，因此有三个红外允许基频，承载五个振动模式。拉曼查 $A_1'$ 与 $E'$，也有三个允许基频、五个模式。两组 $E'$ 在两种光谱中都允许；本分子没有反演中心，不适用红外与拉曼互斥规则。
\end{worked}

\subsection{IR 与 Raman：分别查坐标与二次函数}

第\ref{sec:products}节给出了完整跃迁判据。分析本章的振动基频时，可以把它化成直接查表的方法：

红外电偶极吸收的算符是偶极矩 $\symbf{\mu}$，其三个分量分别按 $x,y,z$ 变换。普通非共振拉曼的算符是极化率张量 $\symbf{\alpha}$，其六个独立分量按 $x^2,y^2,z^2,xy,xz,yz$ 变换。对通常的谐振近似下、同一电子态内的振动基频跃迁，振动基态为全对称，单量子末态具有相应简正坐标 $Q_k$ 的对称性，判据遂简化为
\begin{center}
\small
\begin{tabularx}{\linewidth}{lXX}
\toprule
光谱 & 查特征标表中的基 & 对应的物理条件\\
\midrule
红外 & $x,y,z$ & 至少一个 $(\partial\mu_a/\partial Q_k)_0\ne0$\\
拉曼 & $x^2,y^2,z^2,xy,xz,yz$ & 至少一个 $(\partial\alpha_{ab}/\partial Q_k)_0\ne0$\\
\bottomrule
\end{tabularx}
\end{center}
这里 $a,b$ 是笛卡儿分量指标，下标 0 表示在平衡构型取导数。对称性允许这些导数非零，但导数的实际大小仍由电子结构决定。

\begin{keyidea}[有反演中心时的互斥规则]
偶极矩是 $u$，极化率是 $g$。因此在中心对称分子中，普通基频电偶极红外活性模式只能是 $u$，普通拉曼活性模式只能是 $g$，同一模式不能两者兼有。但不是每个 $u$ 模式都红外活性，也不是每个 $g$ 模式都拉曼活性，仍须查具体基函数。两种谱中恰有相同波数也不等于已证明是同一模式，不能仅凭一次重合否定反演中心。
\end{keyidea}

\subsection{能级模型怎样帮助读谱}

对近似独立的伸缩坐标，谐振子给出
\[
V=\tfrac12 kq^2,\qquad
\widetilde\nu=\frac{1}{2\pi c}\sqrt{\frac{k}{\mu}},\qquad
E_v=hc\widetilde\nu\left(v+\tfrac12\right).
\]
$V$ 是势能，$q$ 为离开平衡键长的位移，$k$ 为力常数，$h$ 为普朗克常数，$c$ 为光速，$\widetilde\nu$ 为波数，$v=0,1,\ldots$ 为振动量子数。双原子分子的约化质量是 $\mu=m_1m_2/(m_1+m_2)$，$m_1,m_2$ 为两原子的质量；多原子简正模式须使用相应有效质量。键越硬，频率通常越高；同位素质量增大，频率通常降低。$v=0$ 仍有零点能 $hc\widetilde\nu/2$。

在谐振势且偶极矩对 $q$ 只保留一次项时，红外振动选律为 $\Delta v=\pm1$，基频吸收为 $0\to1$。非谐振子常写成
\[
\frac{E_v}{hc}=\widetilde\nu_e\left(v+\tfrac12\right)
-\widetilde\nu_e x_e\left(v+\tfrac12\right)^2,
\]
$\widetilde\nu_e$ 为谐振波数，$x_e$ 为无量纲非谐常数。能级间距随 $v$ 增大而减小，故 $0\to2$ 的第一倍频通常略低于基频波数的两倍；非谐性及偶极矩高阶项还可使倍频、组合频获得弱强度，但仍须满足对称性选律。

双原子或线性分子的刚性转子模型为
\[
E_J=hcB J(J+1),\qquad B=\frac{h}{8\pi^2cI},
\qquad I=\mu r_e^2\ \text{（双原子）}.
\]
$J$ 是转动量子数，$I$ 是转动惯量，$r_e$ 是平衡键长，$B$ 是以波数计的转动常数。有永久偶极矩时，纯转动电偶极跃迁满足 $\Delta J=\pm1$；吸收线位于 $2B(J+1)$。例如 $B=1.93\,\mathrm{cm}^{-1}$ 时，$J=1\to2$ 的波数为 $4B=7.72\,\mathrm{cm}^{-1}$。振动谱带在气相还可能带有转动精细结构；本章的群论数带不等于数这些细线。

\section{只分析目标伸缩区}
\label{sec:stretching-applications}
\learning{根据所选内坐标直接建立伸缩表示，用羰基区的理想谱带数区分 fac/mer 异构体，并理解为什么局部伸缩分析不再扣整体平移和转动。}
\sources{D，第40--42页；平面配合物的异构背景见B，第17页，伸缩红外判别见B，第24--25页。}

\subsection{伸缩内坐标为何只数原位键}

若只关心若干条同类键的伸缩，取 $q_j=\Delta r_j$，即第 $j$ 条键的键长变化为基。$q_j$ 是标量：对称操作把一条键映回自身时，拉长仍然是拉长，并不会因为键的空间方向反转而成为压缩；把两条键互换时，则交换两个基函数。因此
\[
\chi_{\mathrm{stretch}}(R)=\text{在操作 $R$ 下映回自身的目标键数}.
\]
这里不乘 $\chi_{xyz}$，因为每条键只贡献一个键长坐标，而不是三个笛卡儿位移。也不再扣平移、转动，因为内坐标对分子的刚体运动本来就不变。前提是所选坐标集合在该点群操作下封闭；若操作把目标键变成未选中的另一类键，就不能单独使用该集合。

伸缩坐标的对称性是精确的分类工具，但真正的简正振动可以混合其他同对称性内坐标。羰基伸缩与较低频的骨架运动通常分离较好，所以用少量 \ce{C-O} 坐标分析羰基区尤其有效。以下以 M 表示中心金属、L 表示相同的非羰基配体、X 表示取代配体，均假定理想骨架，并忽略额外配体取向造成的降对称性。

\subsection{fac 与 mer：先摆坐标，再数羰基}

\begin{worked}[八面体 \ce{ML3(CO)3} 的两种异构体]
\textbf{fac：}三条 \ce{M-CO} 键占据八面体一个面，局部点群为 $C_{3v}$；取 $z$ 沿该面的三重轴。$C_3$ 循环置换三个羰基，每个 $\sigma_v$ 固定一个、交换另两个，故
\[
\chi_{\mathrm{CO}}=(3,0,1)\qquad
\text{列序 }E,2C_3,3\sigma_v.
\]
利用 $A_1:(1,1,1)$、$A_2:(1,1,-1)$、$E:(2,-1,0)$，有
\[
a_{A_1}=\frac{3+3}{6}=1,\quad
a_{A_2}=\frac{3-3}{6}=0,\quad
a_E=\frac{3(2)}6=1.
\]
所以 $\Gamma_{\mathrm{CO}}=A_1+E$。三个伸缩模式中，$E$ 的两个模式简并。因 $z\sim A_1$、$(x,y)\sim E$，通常出现\textbf{两个红外允许的羰基基频}。相应二次函数也存在，故通常有两个拉曼允许基频。

\textbf{mer：}把三个羰基放在 $+z,+y,-y$ 方向，三个 L 放在 $-z,+x,-x$ 方向。主轴为 $C_2(z)$，三个羰基都在 $yz$ 平面，点群为 $C_{2v}$。逐一数原位羰基：
\begin{center}
\small
\begin{tabular}{lrrrr}
\toprule
 & $E$ & $C_2(z)$ & $\sigma(xz)$ & $\sigma(yz)$\\
\midrule
$\chi_{\mathrm{CO}}$ & 3 & 1 & 1 & 3\\
\bottomrule
\end{tabular}
\end{center}
$C_2$ 和 $\sigma(xz)$ 均只固定 $+z$ 羰基；$\sigma(yz)$ 固定全部三个。约化系数依次为
\[
(a_{A_1},a_{A_2},a_{B_1},a_{B_2})
=\tfrac14(8,0,0,4)=(2,0,0,1),
\]
故 $\Gamma_{\mathrm{CO}}=2A_1+B_2$。三个模式均非简并，且分别有 $z$ 或 $y$ 型偶极矩分量，所以通常有\textbf{三个红外允许基频}；它们也都拉曼允许。

结论是理想 fac 的 $2$ 条与 mer 的 $3$ 条羰基红外基频。若把 mer 的分子面改放在 $xz$，会写成 $2A_1+B_1$，物理结论不变。看见三条带时还应先确认它们属于同一种化合物的同一伸缩区，而不是直接把混合物当作一个异构体。
\end{worked}

\subsection{五羰基取代物：简并不等于少一个模式}

\begin{worked}[$C_{4v}$ 的 \ce{M(CO)5X}]
取 X 在 $+z$，一个轴向羰基在 $-z$，四个赤道羰基位于 $\pm x,\pm y$。对 $C_{4v}$ 各类操作，五个 \ce{C-O} 伸缩坐标给出下表。$\sigma_v$ 取 $xz,yz$，$\sigma_d$ 为两条赤道轴之间的平分面。
\begin{center}
\small
\begin{tabular}{lrrrrrl}
\toprule
 & $E$ & $2C_4$ & $C_2$ & $2\sigma_v$ & $2\sigma_d$ & 常用基\\
\midrule
$A_1$ & 1 & 1 & 1 & 1 & 1 & $z;\ x^2+y^2,z^2$\\
$A_2$ & 1 & 1 & 1 & $-1$ & $-1$ & $R_z$\\
$B_1$ & 1 & $-1$ & 1 & 1 & $-1$ & $x^2-y^2$\\
$B_2$ & 1 & $-1$ & 1 & $-1$ & 1 & $xy$\\
$E$ & 2 & 0 & $-2$ & 0 & 0 & $(x,y);\ (xz,yz)$\\
\midrule
$\chi_{\mathrm{CO}}$ & 5 & 1 & 1 & 3 & 1 & \\
\bottomrule
\end{tabular}
\end{center}
例如 $C_4$ 只固定轴向羰基；$\sigma(xz)$ 固定轴向及 $\pm x$ 两个赤道羰基，贡献 3。群阶为 8，逐行内积为
\begin{align*}
a_{A_1}&=(5+2+1+6+2)/8=2,\\
a_{A_2}&=(5+2+1-6-2)/8=0,\\
a_{B_1}&=(5-2+1+6-2)/8=1,\\
a_{B_2}&=(5-2+1-6+2)/8=0,\\
a_E&=[5(2)+1(-2)]/8=1.
\end{align*}
因此
\[
\boxed{\Gamma_{\mathrm{CO}}=2A_1+B_1+E}.
\]
维数 $2+1+2=5$ 正好等于五条目标键。通常有四个不同伸缩基频：两个 $A_1$、一个 $B_1$ 和一个二重简并 $E$。红外只允许 $2A_1+E$，所以是\textbf{三个红外基频}，不是四个；拉曼允许全部表示，所以是\textbf{四个拉曼基频}，不是五个。这里 $B_1$ 是拉曼允许而红外禁戒的伸缩。
\end{worked}

\subsection{与平面异构体相连：只看 M--A 伸缩}

理想平面正方形 \ce{MA2B2} 的 cis 构型取两条 M--A 沿 $+x,+y$。用两键夹角平分线作 $C_2$ 主轴，可得两个伸缩组合为 $A_1+B$，其中 $B$ 的下标取决于另外两轴的约定；它们各与一个坐标分量同对称，因此通常有两个红外允许基频。

trans 构型则可把两条 M--A 放在 $\pm x$，反演交换两键。组合 $q_+=(q_1+q_2)/\sqrt2$ 为 $g$，$q_-=(q_1-q_2)/\sqrt2$ 为 $u$。前者不可能电偶极红外活性，后者按 $x$ 变换，故通常只有一个 M--A 红外允许基频。这一“cis 两条、trans 一条”只数所选 M--A 伸缩区，不是整张红外谱的总峰数，也不把两个 trans 伸缩模式误说成只有一个。

\section{用对称性组织分子轨道}
\label{sec:orbital-applications}
\learning{先按对称性分组，再讨论能量和重叠；能构造四面体氢轨道的 SALC，解释线性分子的 $\sigma/\pi$ 与 $g/u$ 配对，并把方法接到八面体金属 $d$ 轨道。}
\sources{D，第43--53页；$O_h$ 下金属 $d$ 轨道的分组为连接后续配位场理论的应用。}

\subsection{相同对称性只是相互作用的入场条件}

线性组合原子轨道法把分子轨道写成 $\psi=\sum_j c_j\phi_j$，其中 $\phi_j$ 是所选原子轨道，$c_j$ 是展开系数。对称性让我们先把大问题拆成互不混合的子空间：在保持分子对称性的哈密顿算符下，不同不可约表示的轨道之间矩阵元为零；同对称性轨道才可能混合。

不过，“同对称性”不等于“一定强烈成键”。还要有合适的空间重叠、相近的轨道能量。能量差过大时混合很弱；即使对称允许，空间形状也可能使有效重叠很小。因此应先用群论排除不可能的组合，再用能量和重叠判断相互作用大小，而不是凭一个标签画出确定的能级顺序。

SALC 是 symmetry-adapted linear combination，即\textbf{对称性匹配线性组合}：把一组等价配体轨道线性组合成按某个不可约表示变换的函数。它是选择基的方法，不是已经求出的成键分子轨道。实用流程为：选配体轨道并约化其表示；构造并归一化各 SALC；寻找同对称性的中心原子轨道；在各对称性块内组合，检查轨道总数守恒。原始基含几个独立轨道，最终就只能得到几个独立分子轨道，不能把同一个原子轨道重复分配给不同组合。

\subsection{四面体：四个 H 1s 怎样匹配 C 2s、2p}

\begin{worked}[\ce{CH4} 的四个归一化 SALC]
把四个 H 放在正四面体顶点方向
\[
\begin{array}{ll}
1:(1,1,1),&2:(1,-1,-1),\\
3:(-1,1,-1),&4:(-1,-1,1),
\end{array}
\]
这些只表示相对方向，中心 C 在原点。以 $h_1,h_2,h_3,h_4$ 表示四个 H 1s 轨道。它们是球对称的，操作后仍在原位的 H 各贡献 $+1$，故在 $T_d$ 的列序 $E,8C_3,3C_2,6S_4,6\sigma_d$ 下
\[
\chi_{\mathrm H}=(4,1,0,0,2).
\]
取 $A_1:(1,1,1,1,1)$、$T_2:(3,0,-1,-1,1)$，得到
\[
a_{A_1}=\frac{4+8+12}{24}=1,\qquad
a_{T_2}=\frac{12+12}{24}=1.
\]
维数 $1+3=4$ 已用尽全部基，故 $\Gamma_{\mathrm H}=A_1+T_2$。在忽略不同 H 轨道重叠的正交基近似下，可选
\begin{align*}
S_{A_1}&=\tfrac12(h_1+h_2+h_3+h_4),\\
S_x&=\tfrac12(h_1+h_2-h_3-h_4),\\
S_y&=\tfrac12(h_1-h_2+h_3-h_4),\\
S_z&=\tfrac12(h_1-h_2-h_3+h_4).
\end{align*}
每行四个系数的平方和都是 1，任意两行系数内积为 0。$S_{A_1}$ 对所有操作不变；后三个函数整体构成 $T_2$，可按 $x,y,z$ 变换，不能把它们分别当成三个一维不可约表示。若保留轨道间重叠，归一化须改用重叠矩阵，而不再一律是 $1/2$。

C 2s 属于 $A_1$，与 $S_{A_1}$ 配对；C 的 $(2p_x,2p_y,2p_z)$ 属于 $T_2$，分别与 $(S_x,S_y,S_z)$ 配对。四个 H 轨道加四个 C 价轨道，共给八个价分子轨道：一组成键 $a_1$、一组三重简并成键 $t_2$，以及对应的四个反键轨道。八个价电子填满四个成键轨道。C 1s 也是 $A_1$，但能量很低，在常用价轨道图中作为近似不参与成键的芯层轨道处理；不是被对称性禁止混合。
\end{worked}

\subsection{\ce{CO2}：先分 $\sigma/\pi$，再分 $g/u$}

取分子轴为 $z$，O$_A$、O$_B$ 分别位于 $-z,+z$，点群为 $D_{\infty h}$。绕轴旋转不变的轨道属 $\sigma$（严格的不可约表示记号为 $\Sigma$）；一对像 $(p_x,p_y)$ 那样相互变换的轨道属二重简并 $\pi$（$\Pi$）。反演偶、奇分别记作 $g,u$。这里 $g/u$ 是对分子中心反演的性质，不是“成键/反键”的同义词。

\begin{worked}[端氧组合与中心碳轨道怎样配对]
先只看轨道对称性，暂用忽略重叠时的 $1/\sqrt2$ 归一化因子。两个端氧 2s 轨道组成
\[
S_g=\frac{s_A+s_B}{\sqrt2}\sim\sigma_g^+,
\qquad
S_u=\frac{s_A-s_B}{\sqrt2}\sim\sigma_u^+.
\]
上标 $+$ 表示对包含分子轴的镜面反射为对称。p 轨道的相位统一按全局坐标定义，即各 $p_z$ 正瓣都朝 $+z$；反演会同时交换端点并反转 p 轨道符号。因此
\[
Z_g=\frac{p_{z,A}-p_{z,B}}{\sqrt2}\sim\sigma_g^+,
\qquad
Z_u=\frac{p_{z,A}+p_{z,B}}{\sqrt2}\sim\sigma_u^+.
\]
C 2s 为 $\sigma_g^+$，C $2p_z$ 为 $\sigma_u^+$，所以允许的 $\sigma$ 分块分别是 $\{S_g,Z_g,\mathrm C\,2s\}$ 与 $\{S_u,Z_u,\mathrm C\,2p_z\}$。每块有三个基，产生三个轨道；能量相近与否决定混合程度，不能把 C 2s 分别与两套端氧组合各用一次。

对横向分量 $a=x,y$，同样有
\[
P_{u,a}=\frac{p_{a,A}+p_{a,B}}{\sqrt2},\qquad
P_{g,a}=\frac{p_{a,A}-p_{a,B}}{\sqrt2}.
\]
$(P_{u,x},P_{u,y})$ 构成 $\pi_u$，可与 C 的 $(2p_x,2p_y)\sim\pi_u$ 混合，产生一组成键和一组反键 $\pi_u$，各自二重简并。$(P_{g,x},P_{g,y})$ 构成 $\pi_g$，在 C 2s、2p 价轨道基中没有同对称性的中心轨道，因而成为主要位于氧上的\textbf{非键 $\pi_g$ 对}。

价轨道计数也闭合：$\sigma_g^+$ 三个、$\sigma_u^+$ 三个、$\pi_u$ 四个、$\pi_g$ 两个，共 12 个，正是两个 O 各四个和 C 四个价原子轨道的总数。“非键”是这套价轨道模型的描述，不意味着端氧之间的一切相互作用严格为零。
\end{worked}

\subsection{\ce{HCN} 与八面体金属：保留能用的标签}

\ce{HCN} 也是线性分子，但两端不同，点群是 $C_{\infty v}$，没有反演中心，所以\textbf{不能使用 $g/u$ 标签}。取轴为 $z$，H 1s、C 2s、N 2s、C $2p_z$、N $2p_z$ 都属于 $\sigma^+$，共有五个 $\sigma$ 基；C 与 N 的横向 $(2p_x,2p_y)$ 各贡献一套 $\pi$，合计两套 $\pi$、四个轨道。它们仍按 $\sigma$ 与 $\pi$ 分块，但不再有 $\pi_g/\pi_u$ 的配对限制。注意“四个横向轨道”应写作 $2\Pi$ 的维数，而不是一套二维 $\Pi$；轨道能量仍须进一步计算或借助化学模型判断。

最后把同一思路接到理想 $O_h$ 配合物。金属五个 $d$ 轨道分为
\[
(d_{z^2},d_{x^2-y^2})\sim E_g,\qquad
(d_{xy},d_{xz},d_{yz})\sim T_{2g}.
\]
这是二重与三重简并的\textbf{对称性分组}。六个沿坐标轴指向金属的配体 $\sigma$ 轨道组成 $A_{1g}+E_g+T_{1u}$，于是其中的 $E_g$ 可与金属 $d$ 轨道中的 $E_g$ 相互作用；在仅含配体 $\sigma$ 轨道的模型里，没有匹配 $T_{2g}$ 的 SALC。加入配体 $\pi$ 轨道后，$T_{2g}$ 也可以参与相应相互作用。这解释了为什么讨论配位场时先分 $e_g/t_{2g}$，再谈能量高低：标签来自对称性，具体能量分裂和成键性质还来自配体及其相互作用。
